\begin{statementbox}
\textbf{C084\quad 双曲线焦点三角形的内角平分线}

设双曲线
\[
H:\frac{x^2}{a^2}-\frac{y^2}{b^2}=1\quad(a>0,b>0),
\qquad c=\sqrt{a^2+b^2}>a,
\]
左、右焦点分别为 $F_1(-c,0)$、$F_2(c,0)$。
点 $P(x_0,y_0)$ 在双曲线上，且\textbf{不是左右顶点}，即 $y_0\ne0$。
$\angle F_1PF_2$ 的\textbf{内角平分线}与 $x$ 轴相交于 $M(m,0)$。

\textbf{结论一（轴上截距）：}
\[
\boxed{x_0m=a^2}\qquad\Longleftrightarrow\qquad
\boxed{M\left(\frac{a^2}{x_0},0\right)}.
\]
\begin{center}
\begin{tikzpicture}[font=\small,>=stealth]
\draw[gray!40,rounded corners] (0,0) rectangle (14.2,10);
\node[font=\small\bfseries] at (7.1,9.35) {先看主角：内角平分线 PM，也是 P 处切线};

\begin{scope}[shift={(4.5,3.5)},scale=.58]
\draw[->,gray] (-5.8,0)--(6.2,0) node[right] {$x$};
\draw[->,gray] (0,-2)--(0,6.3) node[above] {$y$};
\draw[teal!75!black,thick,domain=-.65:1.05,samples=90,variable=\t]
 plot ({3*cosh(\t)},{4*sinh(\t)});
\draw[teal!75!black,thick,domain=-.65:1.05,samples=90,variable=\t]
 plot ({-3*cosh(\t)},{4*sinh(\t)});

\coordinate (L) at (-5,0); \coordinate (R) at (5,0);
\coordinate (P) at (3.75,3); \coordinate (M) at (2.4,0);
\draw[blue!75!black,thick] (L)--(P);
\draw[orange!85!black,thick] (P)--(R);
\draw[gray] (L)--(R);
\draw[magenta!80!black,very thick] (M)--(P);
\foreach \q in {L,R,P,M} \fill (\q) circle (3pt);
\node[below left] at (L) {$F_1$}; \node[below right] at (R) {$F_2$};
\node[above left] at (P) {$P$}; \node[below left] at (M) {$M$};

\draw[magenta!80!black,thick] (1.7,-1.5556)--(4.9,5.5556);
\draw[magenta!80!black] (P) ++(-161.075:1.05) arc (-161.075:-114.228:1.05);
\draw[magenta!80!black] (P) ++(-114.228:1.3) arc (-114.228:-67.38:1.3);
\node[fill=white,inner sep=1pt] at (2.15,1.85) {$\alpha$};
\node[fill=white,inner sep=1pt] at (3.9,1.1) {$\alpha$};
\fill (0,0) circle(3pt) node[below left] {$O$};
\fill (3,0) circle(3pt); \draw (3,0)--(4,-1.7) node[right,fill=white,inner sep=1pt] {$A(a,0)$};
\end{scope}
\node[align=left,anchor=west] at (9.2,6.2) {紫线把顶角\\分成相等两份};
\node[align=left,anchor=west] at (9.2,4.3) {$P=(x_0,y_0)$\\$M=(m,0)$};
\node[align=left,anchor=west] at (9.2,2.4) {$\boxed{x_0m=a^2}$\\[6pt]右支：$0<m<a$};
\node at (7.1,.65) {图取 $a=3,b=4,\ P=(15/4,3)$；左支同样成立。};

\end{tikzpicture}
\end{center}

\end{statementbox}
\clearpage
\begin{statementbox}
该式对双曲线\textbf{左、右两支均成立}。其中
\[
\begin{cases}
0<m<a,&x_0>a\quad(\text{右支}),\\
-a<m<0,&x_0<-a\quad(\text{左支}).
\end{cases}
\]

\textbf{结论二（角平分线就是切线）：}
\[
\boxed{PM:\frac{x_0x}{a^2}-\frac{y_0y}{b^2}=1}.
\]
即：\emph{双曲线在非顶点 $P$ 处的切线，平分焦点连线所成的内角
$\angle F_1PF_2$。}

\textbf{补充（内切圆的固定切点）：}
若焦点三角形 $PF_1F_2$ 的内切圆与焦点轴 $F_1F_2$ 相切于 $T$，则
\[
\boxed{T=(a,0)\ \text{(右支)};
\qquad T=(-a,0)\ \text{(左支)}}.
\]
\textbf{务必区分：}切点 $T$ 通常不等于角平分线交点 $M$。

\textbf{独立证明补充：}除角平分线定理、焦半径、切线性质、内切圆和
半角向量法外，还可用\textbf{两直线夹角的正切公式（倒角公式）}直接求 $PM$ 的斜率。
该方法会得到内、外两条角平分线，必须以 $-c<m<c$ 选择内角平分线。

\textbf{适用边界：}顶点 $P=(\pm a,0)$ 时焦点三角形退化，
角平分线交点 $M$ 的上述定义不再适用；但顶点处的切线
$x=\pm a$ 可以另行讨论。
\end{statementbox}
